\documentclass[11pt]{amsart}

\usepackage{amsmath}
\usepackage{amssymb}
\usepackage{amsthm}
%\usepackage{psfig}

\begin{document}
\baselineskip=16pt
\textheight=8.5in
\parindent=0pt 
\def\sk {\hskip .5cm}
\def\skv {\vskip .08cm}
\def\cos {\mbox{cos}}
\def\sin {\mbox{sin}}
\def\tan {\mbox{tan}}
\def\intl{\int\limits}
\def\lm{\lim\limits}
\newcommand{\frc}{\displaystyle\frac}
\def\xbf{{\mathbf x}}
\def\fbf{{\mathbf f}}
\def\gbf{{\mathbf g}}

\def\dbA{{\mathbb A}}
\def\dbB{{\mathbb B}}
\def\dbC{{\mathbb C}}
\def\dbD{{\mathbb D}}
\def\dbE{{\mathbb E}}
\def\dbF{{\mathbb F}}
\def\dbG{{\mathbb G}}
\def\dbH{{\mathbb H}}
\def\dbI{{\mathbb I}}
\def\dbJ{{\mathbb J}}
\def\dbK{{\mathbb K}}
\def\dbL{{\mathbb L}}
\def\dbM{{\mathbb M}}
\def\dbN{{\mathbb N}}
\def\dbO{{\mathbb O}}
\def\dbP{{\mathbb P}}
\def\dbQ{{\mathbb Q}}
\def\dbR{{\mathbb R}}
\def\dbS{{\mathbb S}}
\def\dbT{{\mathbb T}}
\def\dbU{{\mathbb U}}
\def\dbV{{\mathbb V}}
\def\dbW{{\mathbb W}}
\def\dbX{{\mathbb X}}
\def\dbY{{\mathbb Y}}
\def\dbZ{{\mathbb Z}}

\def\la{{\langle}}
\def\ra{{\rangle}}
\def\summ{{\sum\limits}}
\def\eps{{\varepsilon}}
\def\lam{{\lambda}}
\def\tr{{\rm tr}}
\def\Mat{{\rm Mat}}



\bf\centerline{Homework \#1. Due on Thursday, Sep 3rd, 11:59pm on Canvas}\rm
\vskip .1cm

\bf\centerline{Plan for next week (Aug 31, Sep 2):}\rm More examples of vector spaces, subspaces, spans, linear independences, bases and dimension  (sections 1.2-1.5 and parts of 1.6 of the book Linear Algebra by Friedberg, Insel and Spence). 
\skv
{\bf Notations:} $\dbN=\{1,2,\ldots\}$ denotes the natural numbers.
Given a field $F$ and $n,m\in\dbN$, the set of $n\times m$ matrices with entries in $F$ is denoted by $\Mat_{n\times m} (F)$.
The notation $\Mat_n(F)$ is a shortcut for $\Mat_{n\times n}(F)$. 
\skv
\bf\centerline{Problems: }\rm
\skv

{\bf Problem 1:} Let $F$ be a field, and let $A\in \Mat_{n\times m}(F)$ and $B\in \Mat_{m\times k}(F)$ for some $n,m,k\in\dbN$. Prove that $(AB)^T = B^T A^T$ where $X^T$ is the transpose of a matrix $X$. Include all the details.

\skv
\bf{Problem 2: }\rm Let $F$ be a field and $n\in\dbN$. Given $A=(a_{ij})\in \Mat_n(F)$, the {\it trace of $A$},
denoted by $\tr(A)$, is the sum of all diagonal entries of $A$:
$$\tr(A)=\sum_{i=1}^n a_{ii}.$$
\begin{itemize}
\item[(a)] Prove that $\tr(\lam A+\mu B)=\lam\, \tr(A)+\mu\, \tr(B)$ for all $A,B\in \Mat_n(F)$ and $\lam,\mu\in F$.
\item[(b)] Let $\mathfrak{sl}_n(F)=\{A\in \Mat_n(F): \tr(A)=0\}$. Prove that $\mathfrak{sl}_n(F)$
is a subspace of $\Mat_n(F)$ (where $\Mat_n(F)$ is considered as a vector space over $F$ with respect
to the usual matrix addition and scalar multiplication).
\item[(c)] Prove that  $\tr(AB)=\tr(BA)$ for all $A,B\in \Mat_{n\times n}(F)$. 
\end{itemize}
\skv

\bf{Problem 3: }\rm Let $F$ be a field. In both parts below $V=F^2=\{(a,b):a,b\in F\}$ as a set. In each part determine (with proof)
whether $V$ is a vector space with respect to given operations. If your answer is yes, find an explicit formula for the zero element
of $V$ (note that the zero element need not equal $(0,0)$):

\begin{itemize}
\item[(a)] addition is defined in the usual way: $(a,b)+(c,d)=(a+c,b+d)$, and scalar multiplication is defined by $\lam(a,b)=(\lam a,b)$ for all $\lam \in F$ and $(a,b)\in V$;
\item[(b)] addition is defined by $(a,b)+(c,d)=(a+c,b+d+1)$, and scalar multiplication is defined by $\lam(a,b)=(\lam a, \lam b+\lam-1)$
\end{itemize}
\skv

\bf{Problem 4: }\rm Let $V$ be the vector space of all polynomials with real coefficients (that is, all functions $f:\dbR\to\dbR$
of the form $f(x)=\sum\limits_{i=0}^n a_i x^i$ for some $n\in\dbN$ and $a_0,\ldots, a_n\in\dbR$). For each of the following
subsets $W$ of $V$ determine (with proof) if $W$ is a subspace of $V$:
\begin{itemize}
\item[(a)] $W=\{f\in V: f(1)\geq 0 \}$;
\item[(b)] $W=\{f\in V: f(1)=f(2) \}$.
\end{itemize}
\skv
\bf{Problem 5: }\rm Let $F$ be a field, $n\in\dbN$ and $V=F^n=\{(x_1,\ldots, x_n): x_i\in F\mbox{ for each }i\}$. We view $V$ as a vector space with respect to the usual addition and scalar multiplication. Let $W=\{(x_1,\ldots,x_n)\in V: \sum\limits_{i=1}^n x_i=0\}$.
\begin{itemize}
\item[(a)] Prove that $W$ is a subspace of $V$.
\item[(b)] Now assume $F$ is finite. Find $|W|$ (with proof). The answer should be expressed in terms of $|F|$ and $n$. {\bf Note:}
As explained in Lecture~1, basic examples of finite fields are $\dbZ_p$ (integers mod $p$) for any prime $p$, but not every finite field
is of this form (even up to isomorphism).  
\end{itemize}
\skv
\bf{Problem 6: }\rm Let $V$ be a vector space over a field $F$. Let
$W_1$ and $W_2$ be subspaces of $V$.
\begin{itemize}
\item[(a)] Prove that $W_1\cup W_2$ is a subspace of $V$ if and only if $W_1\subseteq W_2$
or $W_2\subseteq W_1$.
\item[(b)] Prove that $W_1+W_2$ is a subspace of $V$, where by definition 
$W_1+W_2=\{w_1+w_2: w_1\in W_1, w_2\in W_2\}$. 
\item[(c)] Prove that $W_1+W_2$ is the smallest subspace of $V$ containing $W_1$ and $W_2$,
that is, if $W$ is any subspace of $V$ containing $W_1$ and $W_2$, then $W$ contains $W_1+W_2$.
\end{itemize}
\skv
\bf{Problem 7: }\rm Let $V$ be a vector space over $\dbC$ (complex numbers).
\begin{itemize}
\item[(a)] Explain why $V$ is also a vector space over $\dbR$ (reals) and $\dbQ$ (rationals) with respect to the same operations.
\end{itemize} 
According to (a), given a subset $W$ of $V$, we can ask whether it is a $\dbC$-subspace (that is, a subspace of $V$
considered as a vector space over $\dbC$), and $\dbR$-subspace or a $\dbQ$-subspace. These are 3 different notions.
For example, it is easier to be an $\dbQ$-subspace than an $\dbR$-subspace since to be a $\dbQ$-subspace $W$
only needs to be closed under addition and multiplications by rationals, while to be a $\dbR$-subspace, it needs
to be also closed under multiplications by all reals.

In parts (b) and (c) $V=\dbC$ (where the vector space addition is the addition in $\dbC$ and the scalar multiplication is
the usual multiplication).

\begin{itemize}
\item[(a)] Prove that the only $\dbC$-subspaces of $V$ are $V$ itself and $\{0\}$.
\item[(b)] Describe all $\dbR$-subspaces of $V$ (and prove that there are no other subspaces).
You may use any results from the Chapter 1 of the book, but state your references clearly.
\end{itemize}

\skv

\bf{Problem 8: }\rm Let $V$ be a vector space over a field $F$ and $S=\{v_1,v_2\}$ a subset of $V$ containing precisely two elements.
Prove that $S$ is linearly dependent $\iff$ $v_1$ is a scalar multiple of $v_2$ or $v_2$ is a scalar multiple of $v_1$.
(By definition, $w$ is a scalar multiple of $v$ if $w=\lam v$ for some $\lam\in F$).
\end{document}



